\documentclass[informe]{practicaitic}
\usepackage{lst-vhdl, lst-shell}


\title{Plantilla informe pràctiques}           % Títol pràctica
\assignatura{Sistemes Digitals}                % Assignatura
\numpract{1}                                   % Número pràctica
\autor{Albert Babí \and Sebastià Vila}         % Autors de l'informe



\begin{document}

\section{Insta\l.lació LaTeX}

A Ubuntu/Debian, cal insta\l.lar els següents paquets:

\begin{itemize}
\item texlive
\item texlive-generic-extra
\item texlive-lang-spanish
\item tex-gyre
\item auctex
\item biblatex
\end{itemize}

Per a treballar amb aquesta plantilla i poder fer servir els logos de
la \texttt{UPC} i grau \texttt{iTIC}. Baixeu-vos de l'OCW el tarfile
\texttt{latexbits-X.X.tar} i feu:

\begin{shell}
$ tar xf latexbits-1.0.tar
$ cd LaTeXbits
$ make
\end{shell}

Això us instal·larà un conjunt de paquets locals de \LaTeX{} que
inclouen l'estil per presentar informes de pràctiques.

Una vegada finalitzada l'insta\l.lació d'aquests paquets podeu usar
aquesta classe de la mateixa manera que usarieu les classes
estàndards de documents.


\section{Introducció a LaTeX}

Aquesta és una plantilla per a escriure informes acadèmics amb
\LaTeX. El seu funcionament es basa en macros. A continuació es
presenten algunes exemples.

\subsection{Estructura dels documents}

Un document escrit amb \LaTeX es pot dividir en capítols, seccions i
subseccions. Per exemple, un informe de pràctiques es podria
estructurar de la següent manera:

\begin{verbatim}
\section{Introducció}
\subsection{Descripció de la pràctica}
\subsection{Objectius}
\section{Estudi previ}
\end{verbatim}

Les seccions i subseccions apareixen automàticament a l'índex.


\subsection{Emfatitzat}

Podem emhatitzar text fent \emph{text emfatitzat}.


\subsection{Llistes i enumeracions}

Es poden crear llistes, que continguin diferents nivells, amb la macro
\texttt{itemize}.

\begin{itemize}
\item Podem crear llistes
\item Amb múltiples elements
\item O que continguin subllistes:
  \begin{itemize}
  \item Amb més elements
  \item Com aquests
  \end{itemize}
\item \dots
\end{itemize}

El mateix passa per a les enumeracions, utilitzant la macro
\texttt{enumerate}:
\begin{enumerate}
\item Primer element
\item Element 2
\item Element 3
  \begin{enumerate}
  \item Primer sub-element
  \item Segon sub-element
  \end{enumerate}
\item \dots
\end{enumerate}


\subsection{Inserció de codi}

Exemple d'introducció de codi \texttt{VHDL}.

\begin{vhdl}
library ieee;
use ieee.std_logic_1164.all;

entity shifter is
  port(x        : in std_logic_vector(3 downto 0);
       y        : out std_logic_vector(3 downto 0);
       ctrl     : in std_logic_vector(1 downto 0));
end shifter;

architecture arch of shifter is
begin
  with ctrl select
    y <= x                    when "00",
         x(2 downto 0) & '0'  when "01",
         '0' & x(3 downto 1)  when "10",
         x(0) & x(3 downto 1) when "11",
         "----"               when others;
end arch;
\end{vhdl}





\subsection{Taules}

Per a treballar amb taules s'utilitza la macro \texttt{tabular}.

\begin{center}
  \begin{tabular}{|c|cc|}
    \hline
    Primera columna & Segona columna & Tercera columna \\ \hline\hline
    1 & 2 & 3 \\ \hline
    4 & 5 & 6 \\ \hline
  \end{tabular}
\end{center}




\subsection{Imatges}

Podem incloure imatges amb \texttt{includegraphics}. Pel tipus de
documents que són els informes de pràctiques, és preferible incloure
les figures entremig del text de forma fixa i sense peu de figura.

\begin{center}
  \includegraphics[width=0.7\textwidth]{esquema.png}
\end{center}



\end{document}
